\subsection{\resultrevision{Comparison Design}}
\begingroup
\color{resultrevision}
The main comparison uses MPSLGP, PSLGP and newly trained/tested NSGPII under common scenario-specific objective bounds and empirical reference fronts. PSLGP denotes the IGD-trained PSLGP-I variant~\cite{xu2026pslgp}. All selected outputs are from actual generation 49. The comparison contains 540 MPSLGP task/run tests, 300 PSLGP retests and 300 new NSGPII tests. MPSLGP and PSLGP deploy one fixed conditioned heuristic across 200 preferences; NSGPII deploys its complete final population Pareto-front set. Test-set seed 968356 and 30 configured simulation replications are shared. The main table recomputes 1,140 run-level indicators from raw outcomes, not published summaries.

MPSLGP shares 1,000 individuals within each two- or three-task group, whereas each independent baseline task has 1,000. MPSLGP and the new NSGPII runs are configured for 50 evaluated generations; the PSLGP retests select generation 49 from historical training configured for 51. Thus selected source generations now match, but training schedules and total population/evaluation budgets are not certified identical. MPSLGP's task-specific Top-N selection also differs from PSLGP's recorded-Best selection. For the main table, exact two-sided Wilcoxon signed-rank tests separately compare MPSLGP with PSLGP and NSGPII using 30 paired run IDs within each target and indicator. Raw $p\leq0.05$ defines significance, with no Holm or other multiple-comparison correction. Symbols follow the paired rank-sum direction; these per-comparison tests do not provide family-wise error control. They describe the recorded configurations, not a budget-matched causal effect of transfer. All seeds remain and invalid outcomes are reported and excluded conditionally.

No-transfer is the completed ablation control, not an additional main baseline. It has the same eight groups, seeds 0--29, one-tree crossover, task quotas, IGD training, surrogate procedure, and test configuration as Full. The verified no-transfer profile disables adaptive transfer and sets its transfer probability and budget to zero. All non-ablation training fields match.

We recompute HV, IGD, and GD for Full and no-transfer using one common set of objective bounds and an empirical reference front for each benchmark scenario, pooled only over these two profiles, their groups, and their seeds. Higher HV and lower distances are preferable. Invalid preferences are reported separately; no-transfer quality scores are conditional on valid preferences, and every training seed is retained. The statistical unit is a paired training seed for a fixed group/task target. Means and sample standard deviations are reported over 30 runs. The 54 target/metric comparisons follow the main experiment's exact two-sided Wilcoxon signed-rank convention: zero differences are omitted, absolute ties receive midranks, and raw $p\leq0.05$ defines significance without multiplicity correction. No-transfer-cell symbols give Full's paired rank-sum direction; bold mean rankings are separate. These per-comparison tests do not provide family-wise error control, and a non-significant difference is not evidence of equivalence.

The completed ablation tests the presence of the transfer mechanism under the recorded configurations. It does not individually identify the effects of preference regions, reward terms, task inheritance, or changes in transfer frequency. A budget- and operator-matched independent PSLGP rerun, a $K=1$ controller, and controls that hold transfer frequency constant remain useful follow-up experiments; their results are not claimed in this paper.
\endgroup